\begin{center}
    \large{% Единое место для указания названия работы
Ультразвуковая локальная навигационная система для привязного беспилотного летательного аппарата
} \\[0.3cm]
    \large\textit{Смирнов Егор Ильич} \\[1cm]
    \textbf{Аннотация}
\end{center}

\vspace{0.5cm}

В работе разработана ультразвуковая локальная навигационная система для
определения координат привязного беспилотного летательного аппарата (БПЛА)
относительно наземной платформы (в том числе подвижной) в условиях
недоступности или ненадёжности глобальных навигационных спутниковых систем.
Расстояния до опорных узлов (якорей) измеряются по времени распространения
ультразвукового сигнала; временная синхронизация передающей и приёмной сторон
выполняется по сверхширокополосному (UWB) радиоканалу, а координаты
восстанавливаются методом трилатерации.

Построена математическая модель системы: оценены максимальная дальность работы,
погрешность определения скорости звука во влажном воздухе, влияние геометрии
якорей (GDOP) и проведён сравнительный анализ аналитического метода трилатерации
и метода Гаусса--Ньютона. Разработана и собрана аппаратная часть: передающий узел
БПЛА (ESP32-S3 + DWM1000 с аппаратной синхронизацией по сигналу EXTTXE), тракт
обработки сигнала якоря (INA128P, TL072CN, детектор огибающей, компаратор LM311P) и
изолированная система питания с защитой от наводок; реализован двухимпульсный метод
самокалибровки времени прихода. Экспериментально подтверждён уверенный приём на дистанции
${\approx}34$~м; расчётная погрешность позиционирования~--- порядка $5$~см.

\newpage

\begin{center}
    \textbf{Abstract}
\end{center}

\vspace{0.5cm}

This work develops an ultrasonic local navigation system that estimates the
position of a tethered unmanned aerial vehicle (UAV) relative to a (possibly
mobile) ground platform under conditions where global navigation satellite
systems are unavailable or unreliable. Distances to reference nodes (anchors)
are measured from the time of flight of an ultrasonic signal; the transmitter
and receiver are time-synchronized over an ultra-wideband (UWB) radio link, and
the position is recovered by trilateration.

A mathematical model of the system is built: the maximum operating range, the
error of the speed-of-sound estimate in humid air, and the influence of anchor
geometry (GDOP) are evaluated, and an analytic trilateration method is compared
with the Gauss--Newton method. The hardware is designed and assembled: the UAV
transmitter node (ESP32-S3 + DWM1000 with hardware EXTTXE synchronization), the
anchor signal-processing chain (INA128P, TL072CN, envelope detector, LM311P
comparator) and an isolated, interference-protected power supply; a two-pulse
time-of-arrival self-calibration method is implemented. Reliable reception at ${\approx}34$~m is
confirmed experimentally; the estimated positioning error is on the order of 5~cm.
